\chapter*{Abstract}

Graphene field-effect transistor (GFET) biosensors functionalised with single-stranded
DNA (ssDNA) probes are a promising platform for the rapid, label-free and highly sensitive
detection of pathogens. The performance of such a device depends critically on the capture
probe: to detect a target reliably across clinical strains, the probe must target a
\emph{conserved} genomic region, hybridise strongly and specifically at the assay
temperature, and remain free of stable self-structure that would compete with hybridisation.
Designing such probes by hand does not scale across many genes and organisms, and offers no
reproducible, quantitative ranking of the candidates.

This dissertation presents a reproducible computational pipeline that designs ssDNA capture
probes from publicly available genomic data. Starting from sequences retrieved from the NCBI
databases for six bacterial marker genes, the pipeline performs automatic length clustering,
multiple sequence alignment with MAFFT, and per-column conservation scoring based on the
Percentage of Pairwise Identity (PPI). Sliding-window candidates are screened with
nearest-neighbour thermodynamics (melting temperature, GC content, hairpin and homodimer free
energies) using primer3, and their secondary structure is evaluated with seqfold, condensed
into a continuous ``No-fold'' score. A layered parameter system resolves thresholds by
gene, species and organism type (bacteria, virus, fungus, protozoan or host), so that
AT-rich or GC-rich targets are judged by criteria appropriate to their own genome rather
than by a single global rule. Prioritised candidates are validated three-dimensionally with
Boltz-2, and an existing experimental probe panel is scored with the same metrics for direct
comparison.

Applied to the six target genes, the pipeline reduced 8455 candidate windows to 2943
that passed thermodynamic and structural screening, from which the highest-quality probes
per gene were validated in 3D. Alignment-free diversity and rarefaction analyses were used to
justify the sampling depth and to characterise each gene. The resulting ranked, biophysically
sound probe candidates provide a solid starting point for GFET biosensor fabrication and for
downstream data-driven modelling, and the framework generalises to new genes and organisms
with minimal configuration.

\paragraph{Keywords} graphene biosensor, GFET, ssDNA probe design, multiple sequence
alignment, sequence conservation, secondary structure, 3D structure prediction, pathogen
detection

\cleardoublepage

\chapter*{Resumo}

Os biossensores baseados em transístores de efeito de campo de grafeno (GFET)
funcionalizados com sondas de DNA de cadeia simples (ssDNA) são uma plataforma promissora
para a deteção rápida, sem marcação e altamente sensível de agentes patogénicos. O desempenho
do dispositivo depende de forma crítica da sonda de captura: para detetar um alvo de forma
fiável nas várias estirpes clínicas, a sonda tem de reconhecer uma região genómica
\emph{conservada}, hibridar de forma forte e específica à temperatura do ensaio e não formar
estrutura secundária estável que compita com a hibridação. Desenhar estas sondas manualmente
não é escalável para muitos genes e organismos e não fornece uma ordenação reprodutível e
quantitativa das candidatas.

Esta dissertação apresenta um pipeline computacional reprodutível que desenha sondas de
captura de ssDNA a partir de dados genómicos públicos. Partindo de sequências obtidas nas
bases de dados do NCBI para seis genes marcadores bacterianos, o pipeline realiza a seleção
automática do agrupamento de comprimento, o alinhamento múltiplo com o MAFFT e a pontuação de
conservação por coluna baseada na Percentagem de Identidade Par-a-Par (PPI). As candidatas,
geradas por janela deslizante, são filtradas com termodinâmica de vizinhos mais próximos
(temperatura de fusão, conteúdo em GC, energias de hairpin e homodímero) através do primer3, e
a sua estrutura secundária é avaliada com o seqfold, condensada num score contínuo
``No-fold''. Um sistema de parâmetros em camadas resolve os limiares por gene, espécie e tipo
de organismo (bactéria, vírus, fungo, protozoário ou hospedeiro), de modo a que alvos ricos
em AT ou em GC sejam avaliados por critérios adequados ao seu próprio genoma, e não por uma
regra global única. As candidatas prioritárias são validadas tridimensionalmente com o
Boltz-2, e um painel experimental de sondas já existente é pontuado com as mesmas métricas
para comparação direta.

Aplicado aos seis genes-alvo, o pipeline reduziu 8455 janelas candidatas a 2943 que passaram
a triagem termodinâmica e estrutural, a partir das quais as sondas de maior qualidade por gene
foram validadas em 3D. Análises de diversidade sem alinhamento e de rarefação foram usadas
para justificar a profundidade de amostragem e para caracterizar cada gene. As candidatas
resultantes, ordenadas e biofisicamente sólidas, constituem um ponto de partida robusto para o
fabrico de biossensores GFET e para modelação orientada por dados, e a estrutura generaliza-se
a novos genes e organismos com configuração mínima.

\paragraph{Palavras-chave} biossensor de grafeno, GFET, desenho de sondas ssDNA, alinhamento
múltiplo de sequências, conservação de sequência, estrutura secundária, previsão de estrutura
3D, deteção de agentes patogénicos

\cleardoublepage
